\documentclass[
10pt, % Main document font size
a4paper, % Paper type, use 'letterpaper' for US Letter paper
oneside, % One page layout (no page indentation)
%twoside, % Two page layout (page indentation for binding and different headers)
headinclude,footinclude, % Extra spacing for the header and footer
BCOR5mm, % Binding correction
]{scrartcl}
\input{stru/structure.tex}
\hyphenation{Fortran hy-phen-ation} 

\title{\normalfont{FTML practical session 11}} 
\author{\spacedlowsmallcaps{}
} 
\date{\today}
\begin{document}
\renewcommand{\sectionmark}[1]{\markright{\spacedlowsmallcaps{#1}}}
\lehead{\mbox{\llap{\small\thepage\kern1em\color{halfgray} \vline}\color{halfgray}\hspace{0.5em}\rightmark\hfil}} 
\pagestyle{scrheadings}
\maketitle 
\setcounter{tocdepth}{1} 


\begin{figure}[htpb]
    \centering
    \includegraphics[width=1\linewidth]{Gradient_descent_n_200_d_100_gaussian}
    \label{fig:risks}
\end{figure}

\tableofcontents

\newpage

\section{\large\color{Blue}Convergence speed of gradient descent}

\subsection{Setting}

    We want to study the speed of convergence of the minimization of a convex function $f$ defined over $ \mathbb{R}^d$, with gradient descent.

\begin{equation}
    \label{eq:gd}
    \forall t, \:\theta_{t+1}\leftarrow \theta_t-\gamma \nabla_{\theta}f(\theta_t)
\end{equation}

where $t$ is the iteration index.

We will study the specific case of linear regression (OLS), but the results
parially generalize to general convex functions. We use the usual objects:

    \begin{itemize}
	\item design matrix $X \in \mathbb{R}^{n, d}$
	\item label vector $y\in \mathbb{R}^n$.
        \item loss function
    \begin{equation}
	f(\theta) = \frac{1}{2n} ||X\theta - y||^2
    \end{equation}
    \item $||.||$ is the usual euclidean norm.
    \end{itemize}



    The gradient and the Hessian write:

    \begin{equation}
	\nabla_{\theta}f(\theta) =  \frac{1}{n} X^T(X\theta-y)
    \end{equation}

    \begin{equation}
        H = \frac{1}{n} X^TX
    \end{equation}


    We note $\theta^*$ the minimizers of $f$. All minimizers verify that

    \begin{equation}
	\nabla_{\theta}f(\theta^*) = 0
    \end{equation}

    or

    \begin{equation}
        \label{eq:hessian_star}
        H\theta^* = \frac{1}{n} X^Ty
    \end{equation}

If $H$ is not invertible, they might be not unique, but all have the same function value $f(\theta^*)$.

    \begin{itemize}
    	\item $H$ is symmetric, positive semi-definite.
	\item $H$ is invertible if and only if its smallest eigenvalue $\mu$ is
        $>0$, in which case $f$ is strongly convex (see section 2.3 in \url{https://github.com/nlehir/FTML/blob/master/lecture_notes/lecture%20notes.pdf} )
    \end{itemize}

    \subsection{Convergence speed of gradient descent}

    We assume that $\mu > 0$, meaning that $H$ is invertible. Let us study the convergence speed of GD towards $\theta^*$ (that exsits and is unique).


    \subsubsection{Step 1}%
    \label{ssub:Step 1}
    
    Show that

    \begin{equation}
        \label{eq:taylor}
	\forall \theta \in \mathbb{R}^d, f(\theta)-f(\theta^*) = \frac{1}{2} (\theta-\theta^*)^TH(\theta-\theta^*)
    \end{equation}

    \subsubsection{Step 2}%
    \label{ssub:Step 2}

    Show that

    \begin{equation}
\forall t\in \mathbb{N},	\theta_t  = \theta_{t-1}-\gamma H(\theta_{t-1}-\theta^*)
    \end{equation}

    \subsubsection{Step 3}%
    \label{ssub:Step 3}
    


    Deduce that:

    \begin{equation}
	\theta_t-\theta^* = (I-\gamma H)(\theta_{t-1}-\theta^*)
    \end{equation}

    and that

    \begin{equation}
	\theta_t-\theta^* = (I-\gamma H)^t(\theta_{0}-\theta^*)
    \end{equation}

    where $\theta_0$ is the initial value of $\theta$.


    \subsubsection{Step 4}%
    \label{ssub:Step 4}
    
    
    We can use two measures of performance of the gradient algorithm. Using the
    previous results, they write:

    \begin{itemize}
        \item Distance to minimizer:

	    \begin{equation}
		||\theta_t-\theta^*||^2 = (\theta_0-\theta^*)^T(I-\gamma H)^{2t}(\theta_{0}-\theta^*)
	    \end{equation}
	\item Convergence in function values:
	    \begin{equation}
		f(\theta_t)-f(\theta^*) = \frac{1}{2} (\theta_0-\theta^*)^T(I-\gamma H)^{2t}H(\theta_{0}-\theta^*)
	    \end{equation}

    \end{itemize}
    

    We introduce the \textbf{{condition number}} $\kappa = \frac{L}{\mu} $ where $L$ is the largest eigenvalue of $H$. By convention, if $\mu=0$, $L=+\infty$. Show that with a good choice of $\gamma$, we obtain an \textbf{exponential convergence}

    \begin{equation}
	||\theta_t-\theta^*||^2\leq \big( 1- \frac{1}{\kappa} \big)^{2t}||\theta_0-\theta^*||^2
    \end{equation}

    We note that

    \begin{equation}
	\big( 1- \frac{1}{\kappa} \big)^{2t}\leq \exp(- \frac{1}{\kappa} )^{2t} = \exp(- \frac{2t}{\kappa} )
    \end{equation}
    
    \subsubsection{Simulation}%
    \label{ssub:Simulation}

    Run a simulation that plots both the upper bound and the convergence speed
    of GD on a least squares problem, like seen on Figure~\ref{fig:gdup}. You can adapt the code from \textbf{tp\_06\_line\_search}.

    \begin{figure}[htpb]
        \centering
        \includegraphics[width=1\textwidth]{Gradient_descent_n_200_d_100_gaussian}
        \caption{Comparison of the upper bound and of the actual convergence
        speed.}
        \label{fig:gdup}
    \end{figure}
    

    \subsubsection{Non strongly convex functions}%
    \label{ssub:Non strictly convex functions}
    

    If $\mu=0$, we do not have an exponential convergence guarantee, but rather
    a convergence rate in $\mathcal{O}( \frac{1}{t})$ (the proof is different).

\end{document}
